%% kapa: exam=1 round=19 year=2008 period=1 form=A subject=civil_law subject_ko=민법 seq=5 number=5
%% source: Q-Net 감정평가사 기출문제 (cst00309) artlSeq=1001962 file=extracted/1008/1001962/2008년 제19회 감정평가사 1차시험 기출문제/1교시.hwp
%% conversion: kordoc 4.24.51 -> tools/gen1.py (auto); review=proofread-sonnet; flags=none
\begin{question}{5}
관습법과 사실인 관습에 관한 설명으로 틀린 것은?(다툼이 있는 경우에는 판례에 의함)
\begin{choices}
  \item 사실인 관습은 사적 자치가 인정되는 분야에서는 법률행위의 해석기준으로서 또는 의사를 보충하는 기능으로서 이를 재판의 자료로 할 수 있다.
  \item 당사자의 의사가 명확하지 않은 경우에 임의규정과 다른 관습이 있으면 그 관습에 따른다.
  \item 관습법이 전체 법질서에 부합하지 않게 되었다면 그러한 관습법은 법적 규범으로서의 효력이 부정된다.
  \item 물권은 관습법에 의하여도 창설될 수 있다.
  \item 사실인 관습은 법령으로서의 효력이 없는 단순한 관행이므로, 당사자의 주장이나 증명이 없으면 법원이 직권에 의하여 판단할 수 없다.
\end{choices}
\end{question}
\answer{5}
